\chapter{Was ist eigentlich \enquote{Virtualisierung}?}
\label{kap:grundlagen}

\begin{aufgabe}[Was ist Virtualisierung?]
Stell dir vor, du könntest mehrere Computer gleichzeitig auf einem einzigen Rechner laufen lassen. Genau das ermöglicht Virtualisierung! Erkläre, was man unter Virtualisierung versteht und welche coolen Möglichkeiten das bietet.
\end{aufgabe}

% TODO Antwort

\begin{aufgabe}[Was macht VirtualBox für den Einsatz im Unterricht besonders?]
VirtualBox ist ein Tool, das Virtualisierung richtig einfach macht. Auch im Unterricht verwenden wir diese Software. Aber warum benutzen wir und viele Leute VirtualBox und nicht ein anderes Programm? Was kann VirtualBox gut, und wo gibt es vielleicht bessere Alternativen?
\end{aufgabe}

% TODO Antwort

\begin{aufgabe}[Kann jeder Computer das?]
Nicht jeder Computer kann einfach so mehrere Betriebssysteme auf einmal laufen lassen. Finde heraus, welche Voraussetzungen dein Computer erfüllen muss, damit VirtualBox reibungslos funktioniert.
\end{aufgabe}

% TODO Antwort

\begin{aufgabe}[Vor- und Nachteile der Virtualisierung mit VirtualBox]
VirtualBox klingt fast zu schön, um wahr zu sein! Aber wie bei allem gibt es auch hier zwei Seiten. Welche Vorteile hat VirtualBox, und wo könnte es für dich oder in einem Unternehmen Probleme geben?
\end{aufgabe}

% TODO Antwort

\begin{aufgabe}[Wie erstelle ich meine eigene virtuelle Maschine?]
Du wirst zum Architekten deiner eigenen virtuellen Maschine! Welche Komponenten sind in einer virtuellen Maschine notwendig und was musst du bei diesen einstellen? Gib auch Beispiele für die Werte welche du zuweisen wirst. Begründe warum die Werte nicht beliebig klein und nicht beliebig groß sein dürfen.
\end{aufgabe}

% TODO Antwort

\begin{aufgabe}[Wozu würdest du Virtualisierung auf einem Desktop einsetzen?]
Nenne und beschreibe drei Anwendungsbeispiele, in denen du Virtualisierung auf einem Desktoprechner einsetzen würdest.
\end{aufgabe}

% TODO Antwort

\begin{aufgabe}[Wozu würdest du Virtualisierung auf einem Server einsetzen?]
Nenne und beschreibe drei Anwendungsbeispiele, in denen du Virtualisierung auf einem Server eines Unternehmens einsetzen würdest.
\end{aufgabe}

% TODO Antwort

\begin{aufgabe}[Welche Vorteile bietet die Virtualisierung in Bezug auf IT-Sicherheit und Ressourcenoptimierung?]
\end{aufgabe}

% TODO Antwort
